\section{Why NOT Trusting LLMs for Contextualized Time Series?}
\label{ap:analysis}

We provide a theoretical analysis of why directly trusting LLMs over serialized
contextualized time series can be problematic. The analysis focuses on two
sources of misalignment: numerical values are first mapped into a discrete token
space, and long temporal evidence is then processed under a finite and unevenly
allocated attention budget. We begin with the tokenization-induced mismatch
between numerical distance and token distance.

Proposition~\ref{prop:token_numeric} shows that numerical closeness is not
necessarily preserved after tokenization. This provides a geometric explanation
for why surface-token similarity may be a poor proxy for numerical similarity.

\begin{proposition}[Token Distance Is Not Numerical Distance]
\label{prop:token_numeric}
Let $\eta:\mathbb{R}\rightarrow\mathcal{V}^*$ be a tokenizer-induced
representation from real values to token sequences. In general, edit distance or
embedding distance between $\eta(x)$ and $\eta(x')$ is not order-preserving with
respect to numerical distance $|x-x'|$.
\end{proposition}

\begin{proof}
Tokenizers operate over character or subword patterns rather than numerical
metric structure. Hence there can exist $x_1,x_2,x_3$ such that
\begin{equation}
    |x_1-x_2| < |x_1-x_3|,
\end{equation}
but
\begin{equation}
    d_{\mathrm{tok}}(\eta(x_1),\eta(x_2))
    >
    d_{\mathrm{tok}}(\eta(x_1),\eta(x_3)).
\end{equation}
For example, under a character-level representation,
$x_1=99.999999$, $x_2=100.0$, and $x_3=98.999999$ satisfy
$|x_1-x_2|=0.000001<1.0=|x_1-x_3|$, while $\eta(99.999999)$ and $\eta(100.0)$ differ in more character positions than $\eta(99.999999)$ and $\eta(98.999999)$ (we can alter number of 9s until this holds). Therefore, token similarity does not generally preserve numerical proximity.
\end{proof}

This mismatch is consistent with empirical findings that LLMs often struggle with
numerical reasoning tasks such as arithmetic, numerical retrieval, and magnitude
comparison, partly because numbers are represented through surface token patterns
rather than continuous magnitudes~\citep{li2025exposing}. Prior work also shows that
number-tokenization choices can substantially affect arithmetic performance,
indicating that numerical reasoning is sensitive to the discrete token
representation itself~\citep{singh2024tokenization}. Similar token-level limitations
appear in exact symbolic counting tasks: recent studies report that LLMs can fail
on simple letter-counting problems, where the model must reason over characters
rather than semantic word meaning~\citep{xu2025llm,fu2024large}.
The widely discussed ``strawberry'' counting example is an informal illustration
of the same issue: a word may be represented as subword tokens rather than as
explicit characters, so exact character-level counting is not directly aligned
with the model's token-level representation. This limitation is not unique to
``strawberry''; for example, prior work reports that an LLM can answer that the
word ``confusable'' contains two occurrences of the letter ``a'', while the
correct answer is one~\citep{xu2025llm}. Such examples suggest that even
simple character-level operations may be unreliable when they are mediated by
subword token representations rather than explicit symbolic access to
characters.

Beyond pairwise distance distortion, tokenization can also obscure the internal
structure of numbers. In decimal notation, each digit has a fixed place value,
whereas subword tokenization may group different spans of digits into variable
symbolic chunks. Proposition \ref{prop:place_value} formalizes this place-value mismatch.


\begin{proposition}[Subword Tokenization Does Not Preserve Place-Value Structure]
\label{prop:place_value}
Let $x \in \mathbb{N}$ be a decimal integer with digit expansion
\begin{equation}
    x = \sum_{k=0}^{m-1} d_k 10^k,
    \qquad d_k \in \{0,\ldots,9\}.
\end{equation}
Let $\eta(x)=(v_1,\ldots,v_J)\in\mathcal{V}^*$ be a tokenizer-induced
subword representation of the decimal string of $x$. In general, there does not
exist a fixed token-level decomposition
\begin{equation}
    x = \sum_{j=1}^{J} \phi(v_j) 10^{\alpha_j},
\end{equation}
where $\phi:\mathcal{V}\rightarrow\mathbb{Z}$ and $\alpha_j$ are tokenizer-independent
place-value exponents determined only by token position $j$.
\end{proposition}

\begin{proof}
Decimal notation has a fixed digit-level positional structure: the contribution
of digit $d_k$ is always $d_k10^k$. Subword tokenization, however, partitions the
decimal string into variable-length segments. That is, a number may be tokenized
as chunks of different lengths depending on surface string patterns, vocabulary
entries, and tokenizer merge rules.

Consider two classical numerical constants whose decimal strings may be tokenized
with different segment lengths, e.g.,
\begin{equation}
\begin{split}
    \eta(3.14159) = (3, .14, 159) \\
    \eta(2.71828) = (2.7, 18, 28)
\end{split}
\end{equation}
where tokens are written by their corresponding string pieces for illustration.
While $3.14159$ and $2.71828$ are familiar approximations of $\pi$ and $e$,
their token-level segmentations may expose different symbolic structures. In the
first case, the token ``3'' covers the integer part, ``.14'' covers the decimal
point together with the $10^{-1}$ and $10^{-2}$ places, and ``159'' covers the
$10^{-3}$, $10^{-4}$, and $10^{-5}$ places. In the second case, the token ``2.7''
covers the integer part, the decimal point, and the $10^{-1}$ place, while
``18'' and ``28'' cover the $10^{-2}$--$10^{-3}$ and $10^{-4}$--$10^{-5}$ places,
respectively. Therefore, the numerical place value associated with a token cannot
be determined solely from its token index $j$.

Hence, unless the tokenizer is explicitly constrained to respect digit-level
place values, the token sequence $\eta(x)$ does not provide a fixed token-level
basis aligned with the decimal expansion of $x$. Numerical magnitude must
therefore be reconstructed indirectly from variable-length symbolic chunks,
rather than being natively represented by the tokenization.
\end{proof}



The previous two propositions characterize representation-level mismatches: tokenization does not preserve either numerical geometry or decimal place-value structure. We next show that this mismatch propagates to the learning objective itself: predicting the next token in symbolic space is not equivalent to predicting the next value in numerical space.



\begin{proposition}[Next-Token Prediction Is Not Next-Value Prediction]
\label{prop:ntp_not_nvp}
Let $(\mathcal{V}^*, d_{\mathrm{tok}})$ be the token-sequence space induced by a
tokenizer $\eta:\mathbb{R}\rightarrow\mathcal{V}^*$, and let
$(\mathbb{R}, d_{\mathrm{num}})$ be the numerical space with
$d_{\mathrm{num}}(x,x')=|x-x'|$. Consider a temporal process
$\{X_t\}_{t\geq 1}$ with $X_t\in\mathbb{R}$ and its tokenized observation
$Z_t=\eta(X_t)$. A next-token predictor minimizes
\begin{equation}
    \mathcal{L}_{\mathrm{tok}}(p)
    =
    \mathbb{E}\left[-\log p(Z_{t+1}\mid Z_{1:t})\right],
\end{equation}
whereas a next-value predictor minimizes
\begin{equation}
    \mathcal{L}_{\mathrm{num}}(g)
    =
    \mathbb{E}\left[d_{\mathrm{num}}(g(X_{1:t}),X_{t+1})^2\right].
\end{equation}
In general, an optimizer of $\mathcal{L}_{\mathrm{tok}}$ is not an optimizer of
$\mathcal{L}_{\mathrm{num}}$.
\end{proposition}

\begin{proof}
The next-token objective is a proper scoring rule over the discrete symbolic
space $\mathcal{V}^*$; its Bayes-optimal solution is the conditional token
distribution
\begin{equation}
    p^\star(z\mid Z_{1:t})
    =
    \mathbb{P}(Z_{t+1}=z\mid Z_{1:t}).
\end{equation}
By contrast, under squared numerical loss, the Bayes-optimal next-value predictor is
\begin{equation}
    g^\star(X_{1:t})
    =
    \mathbb{E}[X_{t+1}\mid X_{1:t}].
\end{equation}
These two Bayes solutions live in different spaces: $p^\star$ is a distribution
over token strings, while $g^\star$ is an element of the numerical space
$\mathbb{R}$.

For the two objectives to be equivalent, the tokenizer-induced representation
would need to preserve the task-relevant numerical structure. In particular,
there must exist a decoding map $\delta:\mathcal{V}^*\rightarrow\mathbb{R}$ such
that token-level optimality implies numerical optimality:
\begin{equation}
    \delta\left(
    \arg\max_{z\in\mathcal{V}^*}
    p^\star(z\mid Z_{1:t})
    \right)
    =
    \mathbb{E}[X_{t+1}\mid X_{1:t}]
\end{equation}
for all conditional distributions of $X_{t+1}$. This condition does not hold in
general. First, Proposition~\ref{prop:token_numeric} shows that token distance is
not order-preserving with respect to numerical distance. Therefore, errors that
are small in token space need not be small in numerical space, and vice versa.
Second, the token objective assigns loss according to symbolic likelihood rather
than numerical deviation. For example, predicting the tokenization of $100.0$
instead of $99.999999$ may incur a large token-level discrepancy despite a small
numerical error, while predicting a token string closer in surface form can be
numerically farther away.

Hence, minimizing next-token prediction risk does not generally minimize
next-value prediction risk. The two coincide only under additional assumptions,
such as a value-sufficient tokenizer and a decoding rule that preserves the
numerical geometry relevant to the task loss.
\end{proof}



So far, the analysis concerns the numerical representation induced by
tokenization. Moreover, contextualized time series introduce another challenge: temporal
signals can be long, sparse, and distributed across many timestamps. Even when
all serialized tokens fit into the nominal context window, relevant temporal
evidence must compete with many other tokens under a fixed attention budget.



\begin{proposition}[Long-Context Temporal Information Dilution]
\label{prop:long_context_dilution}
Let a contextualized time series $\mathcal{X}=(\mathbf{X}_{1:T},\mathcal{C})$
be serialized into a token sequence
\begin{equation}
    \mathbf{z}_{1:L}=\operatorname{Tokenize}(\sigma(\mathcal{X})).
\end{equation}
Assume that a temporal task depends on a set of relevant temporal positions
$\mathcal{R}\subseteq [L]$. In a self-attention layer, the output at query
position $i$ is
\begin{equation}
    \mathbf{h}_i
    =
    \sum_{j=1}^{L}\alpha_{ij}\mathbf{v}_j,
    \quad
    \sum_{j=1}^{L}\alpha_{ij}=1,\quad \alpha_{ij}\geq 0 .
\end{equation}
If the attention mass assigned to the relevant positions is bounded by
\begin{equation}
    A_i(\mathcal{R})=\sum_{j\in\mathcal{R}}\alpha_{ij}\leq \epsilon,
\end{equation}
then the contribution of task-relevant temporal evidence to $\mathbf{h}_i$ is
at most $\epsilon$ in attention mass. Therefore, as $L$ grows, any mechanism that
spreads attention over many irrelevant tokens or concentrates attention on
position-biased sink tokens can dilute the effective use of relevant temporal
information.
\end{proposition}

\begin{proof}
Decompose the attention output into relevant and irrelevant parts:
\begin{equation}
    \mathbf{h}_i
    =
    \sum_{j\in\mathcal{R}}\alpha_{ij}\mathbf{v}_j
    +
    \sum_{j\notin\mathcal{R}}\alpha_{ij}\mathbf{v}_j .
\end{equation}
The total coefficient assigned to the relevant component is exactly
\begin{equation}
    A_i(\mathcal{R})=\sum_{j\in\mathcal{R}}\alpha_{ij}.
\end{equation}
If $A_i(\mathcal{R})\leq \epsilon$, then the relevant component contributes no
more than an $\epsilon$ fraction of the convex combination defining
$\mathbf{h}_i$. In particular, if $\|\mathbf{v}_j\|\leq M$ for all $j$, then
\begin{equation}
    \left\|
    \sum_{j\in\mathcal{R}}\alpha_{ij}\mathbf{v}_j
    \right\|
    \leq
    \sum_{j\in\mathcal{R}}\alpha_{ij}\|\mathbf{v}_j\|
    \leq
    \epsilon M.
\end{equation}
Thus, when relevant temporal positions receive little attention mass, their
effect on the attention output is bounded regardless of their task relevance.
Long serialized time series increase the number of irrelevant or weakly relevant
tokens competing for the same unit attention mass. Moreover, if some tokens act
as attention sinks and absorb disproportionate mass, the remaining mass available
to temporally relevant positions is further reduced. This proves the claim.
\end{proof}


This proposition shows that serialized temporal evidence competes within a fixed
attention budget. In long time series, relevant timestamps may be sparse, located
in the middle of the context, or separated by long lags. Empirically, LLMs often
show degraded retrieval when relevant information is placed in the middle of a
long context, a phenomenon known as ``lost in the middle'' \citep{liu2024lost, laban2025llms}.
Meanwhile, attention-sink studies show that some tokens can attract
disproportionate attention mass independent of semantic relevance \citep{xiao2024efficient, gu2025attention, kang2025see}.
Together, these effects suggest that even when a serialized time series fits
within the nominal context window, its temporally relevant information may not be
uniformly or reliably used.


\begin{corollary}[No Guaranteed Length Extrapolation]
\label{cor:length_extrapolation}
Suppose an LLM is trained on serialized sequences of length at most $L_{\max}$.
For a test-time time series whose serialization length satisfies $L>L_{\max}$,
performance on temporal dependencies requiring positions beyond the training
length is not guaranteed by the training objective alone.
\end{corollary}

\begin{proof}
The next-token training objective constrains the model on the support of the
training distribution, whose sequence lengths are at most $L_{\max}$. For
serialized sequences with $L>L_{\max}$, the model is evaluated outside this
length support. Unless the architecture, positional encoding, or training
procedure imposes an extrapolation-consistent structure, there can exist multiple
extensions of the learned conditional distribution that agree on all training
lengths but behave differently beyond $L_{\max}$. Therefore, training loss on
lengths up to $L_{\max}$ alone does not identify a unique or guaranteed behavior
for longer serialized time series.
\end{proof}

Together, these results suggest that text-only LLM workflows face both
representation-level and process-level limitations for contextualized time
series. Tokenization may distort numerical geometry and place-value structure;
next-token prediction optimizes symbolic likelihood rather than numerical
forecasting or decision risk; and long serialized contexts may dilute or miss
temporally relevant evidence. These limitations motivate time-series-native
agentic workflows that operate on temporal objects through structured analysis
rather than relying solely on serialized text.



\section{Fingerprint Computation for Multimodal Memory}
\label{ap:fingerprint}

This section provides details of the fingerprint computation of the multimodal memory in \method. 



\subsection{Context Embedding}
\label{app:phi}

\paragraph{Textual descriptor $c_\tau$.}
For each task $\tau = (q, \mathcal{X}, \mathcal{Y}, y^\star, \ell)$ we form a single textual descriptor $c_\tau$ by concatenating the framing
fields available in $\mathcal{C}$ that identify \emph{what the task is
asking}, rather than the raw numerical signal. The concrete field set
depends on the benchmark schema. For Context-is-Key (CiK), $c_\tau$
joins the \texttt{background}, \texttt{scenario}, and
\texttt{constraints} fields in order. For TSRBench, $c_\tau$ joins the
subtask identifier, domain, task name, the list of channel names, and
the question text. For TSAIA, $c_\tau$ joins the
\texttt{question\_type}, the prompt, the inline data description, and
either the per-task output constraint (for analysis items) or the
option list (for multiple-choice items). Empty fields are skipped, and
the resulting string is truncated to a fixed character budget (we use
$8000$ characters, well below the encoder's $8192$-token limit) so that
pathological cases such as stringified million-element arrays cannot
overflow the encoder.

\paragraph{Encoder.}
The descriptor is embedded by a frozen sentence-level encoder,
\begin{equation}
    \phi = \Phi(c_\tau) \in \mathbb{R}^{d_\phi}.
\end{equation}
We use OpenAI's \texttt{text-embedding-3-small}, which produces
$L_2$-normalized vectors with $d_\phi = 1536$. Because the outputs are
unit-norm, cosine similarity reduces to a dot product,
$\cos(\phi_q, \phi_m) = \phi_q^\top \phi_m$, which we exploit at
retrieval time. If the encoder rejects an input as too long, the
descriptor is halved and re-submitted, with at most two rounds of
halving before the call is treated as a hard failure.

\paragraph{Storage.}
At bank-building time, the embedding is computed once per record and
stored alongside the trajectory; the bank therefore supports retrieval
without re-embedding the entire history at query time. Records logged
before the embedding mechanism was introduced carry a sentinel
no-embedding mask and are simply skipped during the text retrieval
stage.

\subsection{Series Fingerprint}
\label{app:fingerprint}


\begin{table*}[t]
\centering
\small
\caption{The 20 features of the series fingerprint $\psi$, in
assembly order. ``struct'' features describe the channel pool as a
whole; ``per-ch'' features are computed per channel and aggregated by
mean (the standard deviation is additionally log-scaled);
``mv'' is the single multivariate-coupling feature.}
\label{tab:fingerprint_features}
\begin{tabular}{rllp{0.7\textwidth}}
\toprule
\# & Symbol & Type & Description \\
\midrule
1  & $\mathrm{logLen}$              & struct & $\log_{10}(\mathrm{median}_i\, n_i + 1)$; log-scaled median per-channel finite length. \\
2  & $\mathrm{logChan}$             & struct & $\log_2(N' + 1)$; log-scaled number of non-empty channels. \\
3  & $\mathrm{miss}$                & struct & Fraction of entries in $\mathbf{X}_{1:T}$ that are non-finite (NaN or $\pm\infty$). \\
4  & $\mathrm{irreg}$               & struct & Indicator that timestamps are present and unevenly spaced ($\mathrm{CV}(\Delta t) > 0.01$). \\
5  & $\mathrm{meanZ}$               & --     & Reserved slot; identically $0$ in the current implementation. \\
6  & $\mathrm{stdLog}$              & per-ch & $\log_{10}\!\big(\tfrac{1}{N'}\sum_i \sigma_i + 1\big)$; log-scaled mean per-channel standard deviation. \\
7  & $\overline{\mathrm{IQR}/\sigma}$ & per-ch & Mean per-channel IQR-to-std ratio; a robust dispersion fingerprint. \\
8  & $\overline{\mathrm{skew}}$     & per-ch & Mean per-channel skewness on the $z$-scored signal. \\
9  & $\overline{\mathrm{kurt}}$     & per-ch & Mean per-channel excess kurtosis on the $z$-scored signal. \\
10 & $\overline{\mathrm{slope}}$    & per-ch & Mean per-channel $z$-scored linear-fit slope on a unit-rescaled time axis. \\
11 & $\overline{R^2}$               & per-ch & Mean per-channel coefficient of determination of the linear-trend fit. \\
12 & $\overline{\rho(1)}$           & per-ch & Mean autocorrelation at lag $1$. \\
13 & $\overline{\rho(\ell_2)}$      & per-ch & Mean autocorrelation at lag $\lfloor\sqrt{n_i}\rfloor$. \\
14 & $\overline{\rho(\ell_3)}$      & per-ch & Mean autocorrelation at lag $\lfloor n_i/4\rfloor$. \\
15 & $\overline{\mathrm{fftFreq}}$  & per-ch & Mean dominant non-DC frequency from the per-channel rFFT, in cycles per sample. \\
16 & $\overline{\mathrm{fftPFrac}}$ & per-ch & Mean fraction of non-DC spectral power held by the dominant bin. \\
17 & $\overline{\mathrm{specEnt}}$  & per-ch & Mean normalized Shannon entropy of the non-DC spectrum (in $[0,1]$). \\
18 & $\overline{\mathrm{cp}}$       & per-ch & Mean cumulative-sum change-point rate $Z_i / n_i$. \\
19 & $\overline{|\rho|}$            & mv     & Mean absolute Pearson correlation across distinct channel pairs ($0$ if univariate). \\
20 & $\overline{\mathrm{out}}$      & per-ch & Mean per-channel fraction of points beyond $3\,\mathrm{MAD}_i$ from the channel median. \\
\bottomrule
\end{tabular}
\end{table*}


\paragraph{Setup.}
Let $\mathbf{X}_{1:T} \in \mathbb{R}^{N \times T}$ be the numerical
signal of the task and let $\mathbf{X}^{(i)} \in \mathbb{R}^{T}$ denote
its $i$-th channel. After removing non-finite entries (NaN, $\pm\infty$)
we write the finite values of channel $i$ as
$x^{(i)}_1, \ldots, x^{(i)}_{n_i}$, where $n_i \le T$. Channels with
$n_i = 0$ are dropped, and we let $N'$ be the number of remaining
channels. A small constant $\epsilon = 10^{-12}$ is added to
denominators where indicated to keep computations finite.

A per-channel feature is set to $0$ in two degenerate cases: when
$n_i < 3$ (too short for the statistic to be meaningful) and when the
channel standard deviation $\sigma_i$ vanishes (the channel is
constant). These conventions guarantee that the aggregated fingerprint
is always finite, never NaN.

\paragraph{Per-channel statistical features.}
For each channel with $n_i \ge 3$ and $\sigma_i > 0$,
\begin{align}
\mu_i        &= \tfrac{1}{n_i} \textstyle\sum_t x^{(i)}_t, \\
\sigma_i     &= \sqrt{\tfrac{1}{n_i-1} \textstyle\sum_t (x^{(i)}_t - \mu_i)^2}, \\
\mathrm{IQR}_i &= q_{0.75}(x^{(i)}) - q_{0.25}(x^{(i)}), \\
z^{(i)}_t    &= (x^{(i)}_t - \mu_i)/\sigma_i.
\end{align}
On the $z$-scored signal we read off scale-free higher moments,
\begin{align}
\mathrm{skew}_i &= \tfrac{1}{n_i} \textstyle\sum_t (z^{(i)}_t)^3, \\
\mathrm{kurt}_i &= \tfrac{1}{n_i} \textstyle\sum_t (z^{(i)}_t)^4 - 3,
\end{align}
together with a robust outlier rate via the median absolute deviation (MAD),
\begin{align}
\widetilde{x}_i &= \mathrm{median}(x^{(i)}), \\
\mathrm{MAD}_i  &= \mathrm{median}(|x^{(i)} - \widetilde{x}_i|) + \epsilon, \\
\mathrm{out}_i  &= \tfrac{1}{n_i} \big|\{t : |x^{(i)}_t - \widetilde{x}_i| > 3\,\mathrm{MAD}_i\}\big|.
\end{align}

\paragraph{Per-channel trend features.}
A linear trend is fit on a unit-rescaled time axis so that its slope is
comparable across different-length channels. With $u_t = (t-1)/(n_i-1)$
for $t = 1, \ldots, n_i$, let
\begin{equation}
    (a_i, b_i) = \arg\min_{a, b} \textstyle\sum_t \big(x^{(i)}_t - a\, u_t - b\big)^2.
\end{equation}
We record the $z$-scored slope and the linear-fit $R^2$,
\begin{align}
\mathrm{slope}_i &= a_i / (\sigma_i + \epsilon), \\
R^2_i           &= 1 - \frac{\sum_t (x^{(i)}_t - a_i u_t - b_i)^2}
                              {\sum_t (x^{(i)}_t - \mu_i)^2 + \epsilon}.
\end{align}

\paragraph{Per-channel autocorrelation features.}
Let $c^{(i)}_t = x^{(i)}_t - \mu_i$. The sample autocorrelation at lag
$\ell \in \{1, \ldots, n_i-1\}$ is
\begin{equation}
    \rho^{(i)}(\ell)
    = \frac{\sum_{t=1}^{n_i-\ell} c^{(i)}_t c^{(i)}_{t+\ell}}
           {\sum_t (c^{(i)}_t)^2 + \epsilon}.
\end{equation}
We probe three diagnostically useful lags --- a short lag, a
square-root-scale lag, and a quarter-length lag,
\begin{equation}
    \ell^{(i)}_1 = 1, \quad
    \ell^{(i)}_2 = \lfloor \sqrt{n_i} \rfloor, \quad
    \ell^{(i)}_3 = \lfloor n_i/4 \rfloor,
\end{equation}
and record the three values $\rho^{(i)}(\ell^{(i)}_k)$ for $k=1,2,3$.

\paragraph{Per-channel spectral features.}
For $n_i \ge 8$, we compute the real FFT of the centered signal,
\begin{equation}
    \widehat{S}^{(i)}_k = \big| \mathrm{rFFT}(c^{(i)})_k \big|^2,
    \quad k = 0, 1, \ldots, \lfloor n_i/2 \rfloor,
\end{equation}
with bin frequencies $f_k = k/n_i$ in cycles per sample. Let
$k^\star_i = \arg\max_{k \ge 1} \widehat{S}^{(i)}_k$ be the dominant
non-DC bin. We record the dominant frequency and the fraction of
non-DC power it carries,
\begin{align}
\mathrm{fftFreq}_i  &= f_{k^\star_i}, \\
\mathrm{fftPFrac}_i &= \frac{\widehat{S}^{(i)}_{k^\star_i}}
                            {\sum_{k \ge 1} \widehat{S}^{(i)}_k + \epsilon},
\end{align}
together with the normalized Shannon entropy of the non-DC spectrum,
\begin{align}
p^{(i)}_k          &= \widehat{S}^{(i)}_k \Big/
                       \textstyle\sum_{k' \ge 1} \widehat{S}^{(i)}_{k'}, \\
\mathrm{specEnt}_i &= -\frac{1}{\log_2 K_i}
                       \textstyle\sum_{k \ge 1}
                       p^{(i)}_k \log_2 p^{(i)}_k,
\end{align}
where $K_i = \lfloor n_i/2 \rfloor$ is the number of non-DC bins.
Normalizing by $\log_2 K_i$ keeps $\mathrm{specEnt}_i \in [0,1]$
across series of different length. When $n_i < 8$ all three spectral
features are set to $0$.

\paragraph{Per-channel change-point feature.}
Let $d^{(i)}_t = x^{(i)}_{t+1} - x^{(i)}_t$ for $t = 1, \ldots, n_i-1$,
and define the centered cumulative-sum series
\begin{equation}
    S^{(i)}_t = \textstyle\sum_{s=1}^{t}
                 \big( d^{(i)}_s - \overline{d}^{(i)} \big),
\end{equation}
where $\overline{d}^{(i)}$ is the mean of the differences. Letting
$Z_i$ be the number of sign changes in $\mathrm{sign}(S^{(i)})$, we set
\begin{equation}
    \mathrm{cp}_i = Z_i / n_i.
\end{equation}
A stable trajectory produces a CUSUM that drifts smoothly and crosses
zero rarely, whereas a regime-shifting trajectory produces many
crossings; $\mathrm{cp}_i$ is a cheap proxy for piecewise level change.

\paragraph{Structural and multivariate features.}
The remaining features describe the series as a whole rather than any
single channel. The four structural features are
\begin{align}
\mathrm{logLen}   &= \log_{10}\!\big(\mathrm{median}_i\, n_i + 1\big), \\
\mathrm{logChan}  &= \log_2(N' + 1), \\
\mathrm{miss}     &= 1 - \frac{\sum_i n_i}{N\,T}, \\
\mathrm{irreg}    &= \mathbb{1}\!\left[
                       \frac{\sigma(\Delta\mathbf{t})}{\mu(\Delta\mathbf{t})}
                       > 0.01 \right],
\end{align}
where $\Delta\mathbf{t}$ is the sequence of inter-sample time
differences when timestamps are available, and $\mathrm{irreg} = 0$
otherwise. The single multivariate coupling feature is the mean
absolute Pearson correlation across distinct channel pairs,
\begin{equation}
    \overline{|\rho|}
    = \frac{2}{N''(N''-1)} \textstyle\sum_{i<j} |\rho_{ij}|,
\end{equation}
where $N''$ is the number of non-degenerate channels (truncated to a
common length $n_{\min} = \min_i n_i$ and excluded if zero-variance),
and $\rho_{ij}$ is the Pearson correlation between channels $i$ and
$j$. When $N'' < 2$ or $n_{\min} < 3$, $\overline{|\rho|} = 0$.

\paragraph{Aggregation across channels.}
For per-channel features other than the standard deviation, we average
across the $N'$ valid channels,
$\overline{\mathrm{feature}} = \frac{1}{N'} \sum_i \mathrm{feature}_i$.
The standard deviation is averaged and then log-scaled to prevent a
single outlier channel from dominating the metric,
\begin{equation}
    \mathrm{stdLog} = \log_{10}\!\Big( \tfrac{1}{N'}
                       \textstyle\sum_i \sigma_i + 1 \Big).
\end{equation}

\paragraph{Assembly.}
The fingerprint stacks the features in the fixed order listed in
Table~\ref{tab:fingerprint_features},
\begin{equation}
    \psi = \Psi(\mathbf{X}_{1:T}) \in \mathbb{R}^{20}.
\end{equation}
Slot 5 is reserved for a future level-aware feature and is identically
zero in the current implementation; we keep the slot so that
extensions do not invalidate existing banks.

\paragraph{Stability properties.}
Two properties of $\Psi$ are worth noting. First, $\Psi$ is invariant
under permutations of the channel index: every aggregator above is
symmetric in $i$, so a bank populated under one ordering of channels
can still retrieve a query that uses a different ordering. Second, the
use of robust statistics in the outlier-rate, IQR, and median-based
features means that replacing $O(1)$ values per channel with outliers
shifts $\psi$ by $O(1/n_i)$ on those features rather than by an
unbounded amount; the fingerprint thus changes slowly under local
corruption of the input.



\section{Full Memory Construction and Retrieval}
\label{ap:multimodal_memory}

This appendix gives the complete  memory
mechanism summarized in Section~\ref{sec:memory}: the construction of
the context embedding $\phi$, the explicit specification of the series
fingerprint $\psi$, the bank-level normalization that turns the raw
fingerprint into a retrieval-ready key, and the two-stage procedure
that combines the two modalities at query time.

\subsection{Bank-Level Normalization}
\label{app:znorm}

Because the components of $\psi$ live on very different natural scales
--- $\mathrm{logLen}$ is $O(1)$ but $\overline{\mathrm{kurt}}$ can range
over orders of magnitude --- a naive $L_2$ distance would be dominated
by whichever feature happens to have the largest variance in the bank.
We therefore $z$-score each feature using the bank's own per-feature
mean and standard deviation. Let
\begin{align}
    \boldsymbol{\mu}_\psi &= \frac{1}{|\mathcal{M}|}
        \textstyle\sum_{m \in \mathcal{M}} \psi_m, \\
    \boldsymbol{\sigma}_\psi
        &= \mathrm{std}\big(\{\psi_m\}_{m \in \mathcal{M}}\big),
\end{align}
both computed coordinate-wise. Any coordinate with
$\sigma_{\psi, d} < 10^{-9}$ is treated as constant in the bank and is
excluded from the metric by setting $\sigma_{\psi, d} \gets 1$; this
collapses that coordinate to $0$ post-normalization. The
bank-normalized fingerprint is
\begin{equation}
    \tilde\psi = (\psi - \boldsymbol{\mu}_\psi) \oslash \boldsymbol{\sigma}_\psi,
\end{equation}
where $\oslash$ denotes coordinate-wise division. The scaler is
recomputed lazily on every fresh load of the bank rather than persisted,
so that bank growth is reflected immediately at the next query.

\subsection{Two-Stage Retrieval}
\label{app:retrieval}

Given a query task $\tau_q$ with keys $(\psi_q, \phi_q)$, retrieval
proceeds in two stages that combine the two modalities.

\paragraph{Stage 1 (text-side candidate pool).}
We select an intermediate pool of size $N$ by cosine similarity in the
context-embedding space,
\begin{equation}
    \mathcal{N}(\phi_q)
    = \!\!\mathop{\mathrm{Top\text{-}}N}_{m \in \mathcal{M}^\phi}\
    \phi_q^\top \phi_m,
\end{equation}
where $\mathcal{M}^\phi \subseteq \mathcal{M}$ denotes the subset of
records that carry a valid text embedding. The dot-product form is
exact because both $\phi_q$ and the stored $\phi_m$ are unit-norm. In
our experiments we use $N = 20$.

\paragraph{Stage 2 (signal-side ranking).}
Within the text-side candidate pool, we rank by $L_2$ distance in the
bank-normalized fingerprint space,
\begin{equation}
    \mathrm{Retrieve}(\tau_q;\, k)
    = \!\!\mathop{\mathrm{Top\text{-}}k}_{m \in \mathcal{N}(\phi_q)}\
    -\big\| \tilde\psi_q - \tilde\psi_m \big\|_2.
\end{equation}
$k$ is a retrieval parameter that can be adjusted according to the LLM context window and memory density.

\paragraph{Optional filters.}
Two filters can be applied in either stage. A \emph{family filter}
restricts the candidate pool to records sharing the query's task-family
tag, which is useful for diagnostic settings that isolate the
contribution of within-family retrieval. A \emph{self-exclusion filter}
removes the query's own record, which matters when the bank and the
query set overlap (e.g.\ during cross-validation). Both filters are
off by default in our main experiments.

\paragraph{Complexity.}
Stage 1 costs $O(|\mathcal{M}| \cdot d_\phi)$ in time and is dominated
by a single dense matrix-vector product. Stage 2 costs
$O(N \cdot d_\psi)$. At our scale ($|\mathcal{M}|$ in the low
thousands, $d_\phi = 1536$, $d_\psi = 20$, $N = 20$), the entire
retrieval is sub-millisecond on a CPU and does not require an
approximate-nearest-neighbor index.